\documentclass{article}
\usepackage[UTF8]{ctex}
\usepackage{amsmath,amssymb,amsfonts}
\usepackage{enumitem}
\usepackage{geometry}
\geometry{a4paper, margin=1in}
% === 错题本排版优化（不参与题目标识） ===
\setlength{\parindent}{0pt}
\setlength{\parskip}{0.35em plus 0.10em minus 0.08em}
\setlength{\emergencystretch}{3em}
\setlist{itemsep=0.25em, topsep=0.40em, parsep=0pt}
\allowdisplaybreaks[2]
% === 错题本排版优化结束 ===
\begin{document}
\section*{华东师范大学2024年高等代数真题}
\section*{32．华东师范大学2024年高等代数试卷}

一．共10题，每题5分，共50分。

\subsection{题目1}
（1） 设 \(f(x)=x^{3}+3 x+1\) ，满足同余方程 \(v(x) f^{\prime}(x) \equiv 1(\bmod f(x))\) 且次数最小的多项式 \(v(x)\) 为
\(\_\_\_\_\)．

\subsection{题目2}
（2） 实系数多项式 \(f(x)=2 x^{3}-5 x^{2}+4 x-1\) 及 \(g(x)=4 x^{3}-8 x^{2}+5 x-1\) 的最大公因式为 \(\_\_\_\_\)．

\subsection{题目3}
（3）使得实对称矩阵 \(\left(\begin{array}{ccc}3+t & 2 & 1 \\ 2 & 3 & 2 \\ 1 & 2 & 3-t\end{array}\right)\) 正定的实数 \(t\) 的范围为 \(\_\_\_\_\)．
\(\_\_\_\_\)

\subsection{题目4}
（4）已知方阵 \(A\) 的初等因子为 \(\lambda, \lambda^{2},(\lambda+1)^{3},(\lambda-1)^{2},(\lambda-1)^{2}\) ，则 \(A\) 的极小多项式 \(m_{A}(\lambda)\) 为
\(\_\_\_\_\)。

\subsection{题目5}
（5）考虑标准欧氏空间 \(\mathbb{R}^{4}\) 中的向量
\[
\alpha_{1}=(1,1,1,-1), \alpha_{2}=(1,1,1,1), \alpha_{3}=(-1,1,1,1), \beta=(1,-1,1,1) .
\]
设 \(\gamma\) 与诸 \(\alpha_{i}\) 均正交，则 \(\gamma\) 与 \(\beta\) 的夹角的最小值为 \(\_\_\_\_\)．

\subsection{题目6}
（6）设诸 \(\alpha_{i}\) 同上一小题，使得 \(\left|\alpha_{1}-t \alpha_{2}\right|\) 达到最小的实数 \(t\) 为 \(\_\_\_\_\)．

\subsection{题目7}
（7）已知线性映射 \(\varphi: \mathbb{R}^{4} \rightarrow \mathbb{R}^{4}\) 满足
\[
\varphi(x, y, z, w)=(x+3 y+z+6 w, 2 y+z+4 w, 2 x-z,-x+3 y+2 z+6 w) .
\]
那么核空间 \(\operatorname{Ker} \varphi\) 的维数为 \(\_\_\_\_\)．

\subsection{题目8}
（8）置换 \(\tau^{-1} \sigma \tau\) 等于 \(\_\_\_\_\)，这里 \(\sigma=\left(\begin{array}{ccccccc}1 & 2 & 3 & 4 & 5 & 6 & 7 \\ 7 & 6 & 5 & 4 & 3 & 2 & 1\end{array}\right), \tau=\left(\begin{array}{ccccccc}1 & 2 & 3 & 4 & 5 & 6 & 7 \\ 6 & 4 & 2 & 7 & 5 & 3 & 1\end{array}\right)\) ．

\subsection{题目9}
（9）实对称矩阵 \(A=\left(\begin{array}{ccc}4 & -2 & 0 \\ -2 & 3 & -2 \\ 0 & -2 & 2\end{array}\right)\) 可通过正交相似变换化为对角阵 \(\_\_\_\_\)．
\(\_\_\_\_\)

\subsection{题目10}
（10）\(J_{n}(-1)+J_{n}(-1)^{-1}\) 的行列式值为 \(\_\_\_\_\)．

二．每题20分，共100分。

\subsection{题目11}
1 ．已知 \(C(A)=\left\{X \in M_{3}(\mathbb{C}) \mid A X=X A\right\}\) 是 \(\mathbb{C}\) 上的线性空间，这里
\[
A=\left(\begin{array}{lll}
1 & 1 & 0 \\
0 & 1 & 0 \\
0 & 0 & 1
\end{array}\right)
\]
（1）求 \(C(A)\) 的维数．
（2）考虑线性映射
\[
\varphi: C(A) \rightarrow C(A), X \mapsto A X .
\]
求复合映射 \(\varphi^{n}\) 对应的表示矩阵的 Jordan 典范形 \((n \geq 1)\) 。

\subsection{题目12}
2．设
\[
A=\left(\begin{array}{cccc}
-2 & -9 & 2 & -31 \\
0 & -2 & 1 & -5 \\
1 & 4 & 0 & 12 \\
0 & 1 & -1 & 4
\end{array}\right)
\]
（1）求 \(A\) 的特征值及相应的特征子空间．
（2）求 \(A^{2024}\) 。

\subsection{题目13}
3．设 \(A, B \in M_{n}(\mathbb{K}), k \in \mathbb{K}\) 是非零常数。

（1） 证明：如果 \(A B=k I_{n}\) ，则 \(B A=k I_{n}\) 。
\((2)\) 证明：如果 \(A B=O\) ，则 \(\operatorname{rank} A+\operatorname{rank} B \leq n\) 。

\subsection{题目14}
4．设 \(A \in M_{n}(\mathbb{K}), f(x) \in \mathbb{K}[x], \operatorname{Ker}(f(A))=\left\{X \in \mathbb{K}^{n} \mid f(A) X=0\right\}\) 是 \(f(A)\) 的核空间。

（1） 设 \(f(x), g(x) \in \mathbb{K}[x], d(x)=\operatorname{gcd}(f, g)\) ．证明：
\[
\operatorname{Ker}(d(A))=\operatorname{Ker}(f(A)) \cap \operatorname{Ker}(g(A)) .
\]
（2） 设 \(\chi_{A}(x)\) 是 \(A\) 的特征多项式，满足 \(\chi_{A}(x)=f(x) g(x), \operatorname{gcd}(f, g)=1\) ，证明：
\[
\mathbb{K}^{n}=\operatorname{Ker}(f(A)) \oplus \operatorname{Ker}(g(A)) .
\]

\subsection{题目15}
5．设 \(A \in M_{n}(\mathbb{R}), \operatorname{Im}(A)=\left\{Y \in \mathbb{R}^{n} \mid \exists X \in \mathbb{R}^{n}\right.\) 使得 \(\left.Y=A X\right\}\) ．
（1）证明： \(\operatorname{Im}\left(A A^{\mathrm{T}}\right)=\operatorname{Im}(A)\) ．
（2）证明：如果 \(A^{2}=O\) ，那么
\[
\mathbb{R}^{n}=\operatorname{Im}(A) \oplus \operatorname{Im}\left(A^{\mathrm{T}}\right) \oplus \operatorname{Ker}(B)
\]
此处 \(B=A^{\mathrm{T}} A+A A^{\mathrm{T}}, \operatorname{Ker}(B)=\left\{X \in \mathbb{R}^{n} \mid B X=0\right\}\) 。
\end{document}
